\clearpage
\phantomsection

\addcontentsline{toc}{chapter}{{KẾT LUẬN VÀ HƯỚNG NGHIÊN CỨU TIẾP THEO}}

\chapter*{Kết luận và hướng nghiên cứu tiếp theo}

\section*{1. Tổng kết kết quả}

\begin{table}[H]
    \centering
    \small
    \caption{Kết quả đạt được theo từng mục tiêu đề tài, đối chiếu với tiêu chí kết luận ở §\ref{sec:eval_thresholds}}
    \label{tab:mt_results_summary}
    \begin{tabular}{|c|p{3.8cm}|p{4.0cm}|p{2.3cm}|p{1.6cm}|}
        \hline
        \textbf{MT} & \textbf{Mục tiêu} & \textbf{Kết quả \& tiêu chí} & \textbf{Run tốt nhất} & \textbf{Đánh giá} \\
        \hline
        MT1 & Tác nhân chiến đấu hiệu quả trong Roguelike & clear\_rate 0.666, đạt ngưỡng 0.5; 1.58 kills/episode & run62 & Đạt \\
        \hline
        MT2 & Giảm khó khăn do phần thưởng thưa và bản đồ thay đổi & Curriculum learning cải thiện 85\%, vượt ngưỡng 20\% (2.073 $\rightarrow$ 3.843) & run46, run50 & Đạt \\
        \hline
        MT3 & Xác định kỹ thuật hỗ trợ tác động nhất & Curriculum learning tạo cải thiện đáng kể (+85\%); ICM/RND/LSTM chỉ +16--20\% (cải thiện nhẹ) & run46 vs run35/36/42 & Đạt \\
        \hline
        MT4 & Kiểm tra nút thắt quan sát & Entropy giảm 1.35 nat, vượt ngưỡng 1 nat; clear\_rate 0.666 tiệm cận BFS-assisted (0.846) & run62 & Đạt (một seed) \\
        \hline
    \end{tabular}
\end{table}

\noindent Các kết quả thực nghiệm cho thấy đề tài đã đáp ứng các tiêu chí đặt ra cho MT1--MT4 theo các ngưỡng ở §\ref{sec:eval_thresholds}. Tuy nhiên, do mỗi cấu hình chỉ chạy một seed, các kết luận này được hiểu là bằng chứng thực nghiệm có cơ sở, chưa phải kết luận thống kê có khoảng tin cậy.

Khóa luận đã xây dựng và đánh giá một tác nhân học tăng cường sâu cho môi trường game hành động Roguelike trong Unity. Hệ thống cuối cùng sử dụng CombatAgent với quan sát vector 207 chiều, không gian hành động rời rạc nhiều nhánh $(3,3,3,9)$, quy trình sinh bản đồ bằng PCGNN và cơ chế curriculum learning điều khiển qua tham số difficulty. Bản đồ được sinh ngoại tuyến, chuyển sang định dạng văn bản của game, kiểm tra khả năng đi tới bằng BFS và chia thành Easy, Medium, Hard để phục vụ curriculum learning.

Kết quả thực nghiệm cho thấy curriculum learning là xu hướng nổi bật nhất trong nhóm ablation về kỹ thuật huấn luyện. PPO cơ sở trên nhóm bản đồ khó chỉ đạt phần thưởng tail-50 (trung bình 50 lượt chơi cuối) 2.073, còn run46 với curriculum đạt trung bình trên Hard 3.843 -- cải thiện khoảng 85\%. Khi ghép ICM lên nền curriculum, run50 đạt 4.369 phần thưởng và 1.016 kills/episode. Có thể thấy phần thưởng nội tại có tác dụng bổ trợ rõ rệt khi tác nhân đã được làm quen với các bài học dễ hơn. Biến thể chuyển từ vị trí sinh cố định sang ngẫu nhiên trong run54 đạt 3.925, cho thấy tổ chức độ khó môi trường là yếu tố nền quan trọng. Ở nhóm ablation mở rộng về biểu diễn quan sát, run59 với đặc trưng \emph{path-informed} (quan sát có thông tin đường đi BFS) đem lại mức cải thiện đáng kể nhất; bản cập nhật run61 đạt tỷ lệ dọn phòng cuối run khoảng 0.846, số kẻ địch tiêu diệt cuối run khoảng 1.78 và entropy giảm khỏi vùng bão hòa quanh 4.40 nat xuống khoảng 3.62 ở vùng cuối.

Cuối cùng, run62 kết hợp quan sát lưới cục bộ, curriculum 4 bậc, ICM và huấn luyện dài 40.82M bước. Cấu hình này đạt phần thưởng tail-50 14.913, tỷ lệ dọn phòng tail-50 0.666, số kẻ địch tiêu diệt tail-50 1.582 và entropy tail-50 3.048 -- thấp hơn rõ rệt vùng bão hòa 4.40 nat quan sát ở các ablation 1--3M bước với quan sát 207 chiều. Vùng bão hòa này không phải giới hạn cứng của PPO hay không gian hành động. Vì run62 thay đổi nhiều yếu tố đồng thời, nguyên nhân cụ thể vẫn chưa được cô lập hoàn toàn. Toàn bộ kết quả chưa có error bar vì mỗi cấu hình chỉ chạy một seed; chúng nên được diễn giải như bằng chứng thực nghiệm ban đầu thay vì kết luận thống kê cuối cùng.

\begin{table}[H]
    \centering
    \small
    \caption{Các kết quả nổi bật của khóa luận}
    \label{tab:conclusion_best_results}
    \begin{tabular}{|p{4.0cm}|p{3.5cm}|p{6.5cm}|}
        \hline
        \textbf{Kết quả} & \textbf{Giá trị} & \textbf{Ý nghĩa} \\
        \hline
        PPO cơ sở & 2.073 phần thưởng tail-50 & Mốc so sánh không curriculum. \\
        \hline
        Run46 (curriculum) & 3.843 trung bình trên Hard & Tăng khoảng 85\% so với mốc so sánh trong các lần chạy đã thực hiện. \\
        \hline
        Run50 (curriculum + ICM) & 4.369 trung bình trên Hard & Cải thiện thêm trên nền curriculum, nhưng inverse loss vẫn cao. \\
        \hline
        Run54 (cố định--ngẫu nhiên) & 3.925 trung bình trên Hard & Gần run46, cho thấy chuyển từ vị trí sinh cố định sang ngẫu nhiên là một tinh chỉnh môi trường có tác động nhỏ hơn curriculum. \\
        \hline
        Run55 (hàm thưởng v2) & 0.800 giá trị cực đại của tỷ lệ dọn phòng & Tăng giá trị cực đại của tỷ lệ dọn phòng, nhưng phần thưởng trung bình không so trực tiếp do dùng thang thưởng khác. \\
        \hline
        Run59/run61 (\emph{path-informed}) & tỷ lệ dọn phòng cuối run khoảng 0.846, số kẻ địch tiêu diệt cuối run khoảng 1.78 & Bước đột phá về biểu diễn quan sát; bổ sung tín hiệu tìm đường giúp tác nhân tiếp cận và tiêu diệt kẻ địch ổn định hơn. \\
        \hline
        Run62 (lưới cục bộ, 40.82M bước) & 14.913 phần thưởng tail-50, 0.666 dọn phòng, entropy 3.048 & Thoát khỏi vùng bão hòa entropy; tỷ lệ dọn phòng tiệm cận BFS-assisted mà không mã hóa sẵn cơ chế tìm đường. \\
        \hline
        Vùng bão hòa entropy & 4.40 nat (ablation 1--3M bước); 3.048 nat (run62, 40.82M bước) & Run62 thoát khỏi vùng bão hòa cũ khi kết hợp quan sát lưới cục bộ, curriculum, ICM và 40.82M bước huấn luyện. \\
        \hline
    \end{tabular}
\end{table}

Các thí nghiệm ICM, RND và LSTM đem lại thông tin bổ sung nhưng không làm thay đổi xu hướng chính. Khi chạy độc lập, ICM có dấu hiệu inverse model học yếu trên không gian hành động nhiều nhánh; RND cải thiện nhẹ hơn ICM nhưng không vượt trần; LSTM học chậm và tăng dần nhưng chưa tạo cải thiện đáng kể trong quy mô huấn luyện hiện tại. Khi ghép với curriculum, LSTM không có tác dụng bổ trợ rõ, còn ICM trong run50 giúp tăng phần thưởng và số lần tiêu diệt nhưng inverse loss vẫn cao.

\section*{2. Đóng góp của khóa luận}

Khóa luận đem lại sáu đóng góp chính:

\begin{itemize}
    \item \textbf{Quy trình thực nghiệm hoàn chỉnh cho RL trong game hành động Roguelike.} Quy trình trải dài từ thiết kế quan sát/hành động, sinh bản đồ bằng PCGNN, chuẩn hóa bản đồ sang text asset, kiểm tra khả năng đi tới bằng BFS, đến huấn luyện PPO qua Unity ML-Agents. Quy trình không chỉ chạy trên một bản đồ cố định mà còn mở rộng sang tập bản đồ được phân loại theo độ khó.

    \item \textbf{Kết quả thực nghiệm về curriculum learning trong môi trường chiến đấu sinh theo thủ tục.} Đưa tác nhân qua chuỗi Easy -- Medium -- Hard giúp chính sách học được các kỹ năng cơ bản trước khi đối mặt với nhóm bản đồ khó, đem lại cải thiện lớn hơn so với phần thưởng nội tại hay thay đổi kiến trúc bộ nhớ khi các kỹ thuật này chạy độc lập.

    \item \textbf{Quy trình tuyển chọn bản đồ PCGNN phục vụ curriculum.} Hệ thống bổ sung bộ chuyển đổi đặt \texttt{P}/\texttt{E}, kiểm tra đường đi, tính các chỉ số tỷ lệ tường, độ dài đường đi, tỷ lệ ngõ cụt, A* difficulty và chia bản đồ theo percentile. Với lô 14$\times$14 gồm 1000 bản đồ, hệ thống tạo được 1000 bản đồ duy nhất, 100\% bản đồ có đường đi hợp lệ và chia thành 50 Easy, 50 Medium, 900 Hard để người thiết kế lựa chọn.

    \item \textbf{Các kết quả ablation có giá trị.} ICM khi chạy độc lập hoạt động yếu trên không gian hành động rời rạc nhiều nhánh của bài toán, thể hiện qua inverse loss quanh 4.36 so với mức ngẫu nhiên lý thuyết 5.49. Khi ghép với curriculum, run50 cải thiện phần thưởng Hard lên 4.369 nhưng inverse loss vẫn ở 4.337, cho thấy tín hiệu curiosity có thể hỗ trợ hành vi nhưng chưa giải quyết triệt để vấn đề biểu diễn.

    \item \textbf{Phát hiện và thoát khỏi vùng bão hòa entropy.} Trong nhóm ablation 207 chiều, entropy duy trì quanh 4.40 nat bất kể thay đổi ICM, RND, LSTM, curriculum, giảm phương sai vị trí sinh hay thiết kế lại phần thưởng. Cấu hình run62 thoát khỏi vùng bão hòa này khi kết hợp quan sát lưới cục bộ, curriculum 4 bậc, ICM và 40.82M bước huấn luyện: entropy giảm xuống 3.048 nat và duy trì bền vững. Như vậy, vấn đề nằm ở biểu diễn quan sát và quy mô huấn luyện chứ không phải ở kiến trúc PPO hay không gian hành động.

    \item \textbf{Ablation mở rộng về biểu diễn quan sát qua run59.} Khi bổ sung đặc trưng \emph{path-informed} (quan sát có thông tin đường đi BFS), tác nhân đạt mức cải thiện lớn nhất về hành vi: tỷ lệ dọn phòng, số kẻ địch tiêu diệt và mức độ quyết đoán của chính sách đều tốt hơn rõ rệt. Có thể kết luận biểu diễn không gian là yếu tố quyết định trong môi trường chiến đấu sinh theo thủ tục, mở ra hướng kết hợp RL với các đặc trưng điều hướng có cấu trúc thay vì chỉ thay đổi thuật toán huấn luyện.
\end{itemize}

\section*{3. Hạn chế}

Khóa luận còn một số hạn chế cần ghi nhận:

\begin{itemize}
    \item \textbf{Mỗi lần chạy mới chỉ dùng một seed.} Vì vậy, các kết luận về chênh lệch nhỏ cần được hiểu như xu hướng thực nghiệm, chưa phải kết quả thống kê có error bar. Các phát hiện có chênh lệch lớn (như curriculum tăng khoảng 85\%) đáng tin cậy hơn so với các chênh lệch nhỏ; tuy nhiên vẫn cần nhiều seed để khẳng định độ ổn định.

    \item \textbf{PCGNN chạy ngoại tuyến kèm bước tuyển chọn thủ công.} Cách làm này giúp ổn định huấn luyện và giảm lỗi tích hợp, nhưng chưa phải cơ chế sinh nội dung thích ứng theo thời gian thực với độ khó bản đồ điều chỉnh trực tiếp theo năng lực hiện tại của tác nhân. Khóa luận cũng chưa có đánh giá tách biệt đủ sâu về khả năng khái quát hóa trên bản đồ chưa từng thấy. Ví dụ kiểm thử cần thiết: so sánh tác nhân huấn luyện trên bản đồ thủ công với tác nhân huấn luyện bằng curriculum PCGNN, rồi kiểm tra cả hai trên tập bản đồ PCGNN giữ riêng.

    \item \textbf{Quy mô huấn luyện chưa đồng nhất giữa các nhóm.} Nhóm không dùng curriculum chủ yếu dừng ở 1.0M bước, nhóm curriculum chạy đến 3.0M bước, còn run62 đạt 40.82M bước. Một số so sánh do đó phản ánh hiệu quả của cấu hình hoàn chỉnh hơn là ablation công bằng tuyệt đối theo cùng số bước.

    \item \textbf{Điều hướng trong bản đồ sinh theo thủ tục là nút thắt còn lại.} Run59 cho thấy bài toán này chưa được giải quyết hoàn toàn bằng quan sát cục bộ 207 chiều. Khi bổ sung hướng bước tiếp theo theo BFS và khoảng cách đường đi ngắn nhất vào tín hiệu học, hiệu năng tăng rõ rệt. Có thể xem run59 như ablation mở rộng về biểu diễn quan sát: yếu tố tìm đường có tác động lớn tới khả năng tiếp cận và tiêu diệt kẻ địch.

    \item \textbf{Nguyên nhân chính xác của vùng bão hòa entropy chưa được cô lập.} Run62 thoát khỏi vùng bão hòa bằng cách kết hợp nhiều thay đổi đồng thời (lưới cục bộ, curriculum 4 bậc, ICM, hyperparameter điều chỉnh và số bước huấn luyện lớn hơn), nhưng chưa có ablation đơn biến nào xác định yếu tố nào là nguyên nhân chính. Một số hướng như giảm số hướng aim, gộp nhánh combat--aim, thay đổi hệ số entropy, factorized policy hay khởi tạo bằng imitation learning vẫn chưa được kiểm tra.
\end{itemize}

\section*{4. Hướng nghiên cứu tiếp theo}

Trên cơ sở các phát hiện thực nghiệm, khóa luận đề xuất bảy hướng phát triển tiếp theo:

\begin{itemize}
    \item \textbf{Thiết kế lại không gian hành động.} Không gian $(3,3,3,9)$ tuy đủ để biểu diễn điều khiển game nhưng làm entropy cao và khiến chính sách khó trở nên quyết đoán. Có thể gộp nhánh combat/aim, giảm số hướng aim, hoặc tăng cường action mask theo ngữ cảnh.

    \item \textbf{Chạy lại các lần chạy quan trọng với nhiều seed.} Tối thiểu nên chạy run34, run46, run50, run54 và run55 với ba seed để có trung bình, độ lệch chuẩn và khoảng tin cậy, giúp các kết luận về curriculum, ICM trên nền curriculum và hàm thưởng v2 có sức thuyết phục thống kê cao hơn.

    \item \textbf{Tiếp tục khai thác quan sát lưới cục bộ.} Run62 cho thấy khi tăng số bước huấn luyện lên 40.82M, lưới cục bộ có thể đạt tỷ lệ dọn phòng tiệm cận BFS-assisted mà không cần mã hóa sẵn cơ chế tìm đường. Bước tiếp theo là tăng độ phân giải lưới, kết hợp với quan sát pixel hoặc đặc trưng đồ thị để kiểm tra khả năng khái quát hóa sang các bản đồ chưa từng thấy.

    \item \textbf{Khởi tạo bằng imitation learning hoặc behavior cloning.} Với một tập minh họa nhỏ từ người chơi, tác nhân có thể bắt đầu từ chính sách biết di chuyển, né và tấn công cơ bản, rồi dùng PPO/curriculum để tinh chỉnh trên nhóm bản đồ khó.

    \item \textbf{Chuyển PCGNN từ ngoại tuyến sang curriculum thích ứng.} Khi tác nhân đạt hiệu năng cao ở một nhóm độ khó, hệ thống có thể sinh thêm bản đồ mới quanh vùng khó hơn, kiểm tra bằng chỉ số định lượng và đưa vào tập huấn luyện. Hướng này gần với sinh nội dung thủ tục thích ứng (adaptive PCG) và có thể giảm hiện tượng tác nhân học quá khớp vào tập bản đồ cố định.

    \item \textbf{Tăng quy mô huấn luyện có hệ thống.} Run62 đạt 40.82M bước và vẫn cho thấy xu hướng cải thiện; điểm bão hòa chưa xác định được. Huấn luyện thêm kết hợp với đánh giá định kỳ trên tập bản đồ giữ riêng sẽ giúp xác định khi nào nên dừng hoặc chuyển sang tinh chỉnh kiến trúc.

    \item \textbf{Mở rộng sang huấn luyện multi-agent.} Môi trường TopDown Engine hỗ trợ nhiều nhân vật và nhiều kẻ địch. Các kịch bản hợp tác hoặc đối kháng là hướng tự nhiên để kiểm tra khả năng phối hợp, chia mục tiêu và thích nghi chiến thuật của tác nhân.
\end{itemize}

Mã nguồn và tài liệu liên quan được cập nhật tại:
\url{https://github.com/tranmanh2004}
